\documentclass[11pt]{article}
\usepackage{amsmath,amssymb,amsthm,mathtools,enumitem,cleveref,xr}

\externaldocument[uomsrc-]{source_internal_control_factorization_note}
\externaldocument[uommain-]{main}
\externaldocument[stack-]{odd_characterization_transport_stack_note}

\newtheorem{theorem}{Theorem}
\newtheorem{lemma}[theorem]{Lemma}
\newtheorem{proposition}[theorem]{Proposition}
\newtheorem{corollary}[theorem]{Corollary}
\newtheorem{remark}[theorem]{Remark}
\theoremstyle{definition}
\newtheorem{definition}[theorem]{Definition}

\newcommand{\Hacc}{\mathcal H_{\mathrm{acc}}}
\newcommand{\SPhys}{\mathcal S_{\mathrm{phys}}}

\title{Post-HLS $\mathbb Z_2$ Rigidity Over Accepted-Record History}
\date{2026-03-19}

\begin{document}
\maketitle

\section*{Purpose}

This note isolates the downstream post-Main lane that begins only after the UOM\(\to\)Main export
boundary.
It is not part of the weak Main-facing characterization interface and it does not modify the base
\(\mathbf{2T}_{\mathrm{char}}\) contract.
Its purpose is to package the remaining broad sign-obstruction evidence and the post-HLS
\(\mathbb Z_2\)-rigidity target over accepted-record history.

The Main-facing base used here is imported from
the UOM Reduction and Pulse-Rigid Ambiguity Supplement note:
\begin{enumerate}[leftmargin=2em]
\item the accepted-record reduction theorem
  \textup{(\cref{uomsrc-thm:accepted-record-reduction})};
\item the reduced welded-image ambiguity theorem
  \textup{(\cref{uomsrc-prop:reduced-welded-ambiguity-trivial})};
\item the pulse-rigid ambiguity package
  \(
  \mathbf{2T}_{\mathrm{char\!-\!amb}}
  \)
  and its UOM export theorem
  \textup{(\cref{uomsrc-def:2T-char-amb,uomsrc-prop:char-ambiguity-supplement-export})};
\item the reduced physical null-ideal descent theorem
  \textup{(\cref{uomsrc-prop:full-ambiguity-narrowed-bridge})};
\item the exact remaining full-sector identification boundary
  \textup{(\cref{uomsrc-rem:full-sector-gap})}.
\end{enumerate}
The bounded completion-side lane from the Split-Controlled Completion Lift note is an additional
downstream input, not part of the upstream Main-facing supplement; status clarifications for that
optional lane are archived in
\texttt{UOM/post\_hls\_z2\_route\_status\_memo.md}.

\section{Broad Sign Obstruction}

\begin{proposition}[Seam/filler sign obstruction on the broad realization class]
\label{prop:broad-tau-obstruction}
If one identifies the unreduced welded control with the full admissible welded odd trace-channel
source class \(\tau\), then
\[
\tau^{\mathrm{bridge},y}_{\Lambda_k}
=
-\,\tau^{\mathrm{diag},y}_{\Lambda_k}
\]
is compatible with identical accepted-record history while still flipping the unreduced projected
trace-channel sign data.
\end{proposition}

\begin{proof}
\noindent\textbf{Step 1 [definition-level].}
By the Active Response Transport note \textup{(\cref{stack-def:witness-seam-realization})}, there
are live accepted-tick realizations with
\[
\tau^{\mathrm{bridge},y}_{\Lambda_k}
=
-\,\tau^{\mathrm{diag},y}_{\Lambda_k}
\]
and identical accepted-record payload.

\noindent\textbf{Step 2 [definition-level plus theorem-backed already].}
By the UOM main manuscript \textup{(paragraph~\ref{uommain-par:uom-main-081})}, the linearized Gaussian
receiver response is
\[
\delta C=\mathsf T_\Lambda \tau.
\]
By the Active Response Transport note \textup{(\cref{stack-prop:seam-scalar-signs})}, the projected
trace-channel data satisfy
\[
\Pi_{\mathrm{phys},\Lambda_k}\mathsf T_{\Lambda_k}\!\big[\tau^{\mathrm{bridge},y}_{\Lambda_k}\big]
=
-\,\Pi_{\mathrm{phys},\Lambda_k}\mathsf T_{\Lambda_k}\!\big[\tau^{\mathrm{diag},y}_{\Lambda_k}\big].
\]

\noindent\textbf{Step 3 [conclusion].}
Therefore equal accepted-record history does not determine the full broad unreduced realization
data. This is not a Main-facing export theorem; it is source-backed evidence for a surviving hidden
sign involution on the broad realization class.
\end{proof}

\begin{remark}[Physical meaning of the obstruction]
\label{rem:broad-tau-obstruction-physical}
The obstruction is naturally read as an interleaved history effect rather than a same-aeon receiver
ambiguity.
Equal accepted-record history fixes the current reduced receiver object, but it does not fix the
upstairs odd completion datum that survives across the split and can reappear only through later
downstream readout.
That is why the present proposition is evidence for a post-HLS \(\mathbb Z_2\) lift over
accepted-record history rather than evidence for a same-aeon factorization theorem on the broad
source class.
\end{remark}

\section{Post-HLS Target Package}

\begin{definition}[Physical-quotient post-HLS odd package]
\label{def:post-hls-package}
Fix an accepted welded branch \(\beta\).
A \emph{physical-quotient post-HLS odd package} over \(\beta\) consists of the data actually used by
the downstream Main post-HLS lane:
\begin{enumerate}[leftmargin=2em,label=(Z\arabic*)]
\item the physical odd derivation descending from the UOM-fed Main input;
\item its UV spinorial realization \(Q\);
\item the UV dilatation action from Subprogram~4a;
\item the \(S\)-partner structure generated from \(K_\mu\) in Subprogram~4b;
\item the superconformal closure data used in Subprograms~4a--4c, up to the level actually invoked.
\end{enumerate}
Write \(\SPhys(\beta)\) for the set of such physical-quotient packages on a fixed Main background
stratum, and \(\Hacc(\beta)\) for the accepted-record histories on that same stratum.
\end{definition}

\begin{remark}[Retargeted downstream theorem]
\label{rem:z2-lift-target}
The active downstream theorem target is no longer:
\begin{enumerate}[leftmargin=2em]
\item a secondary compression from weak odd-lane data back to full accepted-record history, or
\item an unreduced compiled-control factorization theorem on the broad welded realization class.
\end{enumerate}
The correct downstream target is instead the post-HLS \(\mathbb Z_2\)-rigidity problem over
accepted-record history:

\medskip
\noindent\emph{construct a surjective map}
\[
\pi_{\mathrm{lift}}:\SPhys(\beta)\to\Hacc(\beta)
\]
\emph{whose fibers are singleton or free orbits of a canonical involution \(\iota\), with}
\(\iota\) \emph{acting trivially on accepted-record history and on the reduced welded receiver image.}
\medskip

The broad sign obstruction of \cref{prop:broad-tau-obstruction} supplies the candidate hidden
involution upstairs.
The imported Main-facing base from the UOM Reduction and Pulse-Rigid Ambiguity Supplement note
\textup{(\cref{uomsrc-thm:accepted-record-reduction,uomsrc-def:2T-char-amb,uomsrc-prop:full-ambiguity-narrowed-bridge,uomsrc-rem:full-sector-gap})}
provides the theorem-visible reduction result, the narrow ambiguity package, the reduced null-ideal
descent theorem, and the exact remaining full-sector boundary.
What remains genuinely new is the post-HLS rigidity step excluding any hidden modulus beyond that
candidate sign sheet.
\end{remark}

\section{Dependency Boundary}

\begin{remark}[What belongs upstream and what belongs here]
The pulse-rigid ambiguity supplement belongs upstream of Main as an imported UOM-specific theorem on
one realization of the weak interface.
The present note starts strictly downstream of that boundary.
For the reduced theorem-facing route, the relevant transfer step is already closed by
\textup{\cref{uomsrc-prop:full-ambiguity-narrowed-bridge}}.
If one wants the literal full physical-observable formulation of the post-HLS package, the
remaining upstream task is instead the broader AQFT/HLS identification boundary isolated in
\textup{\cref{uomsrc-rem:full-sector-gap}}.
The present \(\mathbb Z_2\)-rigidity problem is therefore already the correct next downstream target
on the reduced route, while the full-sector formulation still waits on that one later identification
theorem.
\end{remark}

\begin{remark}[Completion-side lane stays separate]
The bounded completion-side package of the Split-Controlled Completion Lift note is a downstream
proving lane for this rigidity problem.
Its explicit split-data blocker is not part of the weak Main-facing ambiguity supplement and should
not be folded back into \(\mathbf{2T}_{\mathrm{char}}\).
The archived status boundary for that optional lane is recorded in
\texttt{UOM/post\_hls\_z2\_route\_status\_memo.md}.
\end{remark}

\end{document}
